\section{Related Work}

\para{Distributional Reinforcement Learning} The foundational work of \citet{bellemare2017distributional} introduced the \cfiftyone algorithm, which models the return distribution using a discrete categorical distribution. This approach was later refined by \citet{dabney2018distributional} with \qrdqn, which uses quantile regression to learn a fixed set of quantiles, and \citet{dabney2018implicit} with \iqn, which learns an implicit representation of the quantile function. While these methods have set the standard for \distrl, they all impose specific structural assumptions on the distribution (either supports or quantiles). Our work, in contrast, uses unconstrained particles to represent the distribution, providing more flexibility in capturing the underlying return dynamics.

\para{Moment Matching and MMD in RL} The use of Maximum Mean Discrepancy (\mmd) for distribution matching in RL was pioneered by \citet{nguyen2020distributional} with \mmddqn. They proved the contraction property of the distributional Bellman operator under the \mmd metric and demonstrated its effectiveness on Atari benchmarks. Building on this, our proposed \immq utilizes \mmd but incorporates the inductive bootstrapping framework from \imm \cite{zhou2025inductive}. This inductive approach provides a more stable training signal and a principled way to incorporate a mean-matching regularizer, which we show is crucial for effective learning in continuous control tasks.

\para{Inductive Moment Matching} Originally proposed for generative modeling, \imm \cite{zhou2025inductive} provides a single-stage training procedure that learns a one-step sampler by matching moments through an inductive objective. It circumvents the need for multi-stage distillation or complex GAN-like training by utilizing marginal-preserving interpolants and \mmd. Our work is the first to adapt this "inductive" philosophy to the bootstrapping process in Q-learning, specifically for the purpose of matching the return distribution to the shifted target distribution in a single, stable training loop.
